\documentclass[11pt]{article}
\usepackage[a4paper,margin=26mm]{geometry}
\usepackage[T1]{fontenc}
\usepackage{lmodern}
\usepackage{microtype}
\usepackage{amsmath,amssymb,amsthm,mathtools}
\usepackage{booktabs}
\usepackage{enumitem}
\usepackage{xcolor}
\usepackage[colorlinks=true,linkcolor=blue!55!black,citecolor=blue!55!black,urlcolor=blue!55!black]{hyperref}
\hypersetup{pdftitle={The Exact Rank Floor for Point-Span Tangent Absorption in Arbitrary Characteristic},pdfauthor={Lluis Eriksson},pdfsubject={Algebraic geometry; fat points; finite projection; tangent spaces},pdfkeywords={fat points, symbolic square, Frobenius, finite projection, tangent absorption, interpolation}}

\newtheorem{theorem}{Theorem}[section]
\newtheorem{proposition}[theorem]{Proposition}
\newtheorem{lemma}[theorem]{Lemma}
\newtheorem{corollary}[theorem]{Corollary}
\newtheorem{remark}[theorem]{Remark}
\newtheorem{definition}[theorem]{Definition}
\newcommand{\PP}{\mathbf P}
\newcommand{\CC}{\mathbf C}
\newcommand{\kk}{\mathbf k}
\newcommand{\NN}{\mathbf N}
\newcommand{\cI}{\mathcal I}
\newcommand{\cO}{\mathcal O}
\newcommand{\Span}{\operatorname{Span}}
\newcommand{\rk}{\operatorname{rk}}
\newcommand{\Gr}{\operatorname{Gr}}
\newcommand{\BZ}{B_d(m)}

\title{The Exact Rank Floor for Point-Span Tangent Absorption\\
in Arbitrary Characteristic}
\author{Lluis Eriksson\\\small Independent researcher}
\date{August 24, 2026\quad (working version 0.7)}

\begin{document}
\maketitle

\begin{abstract}
Let $X$ be a smooth projective integral variety of dimension $d\geq1$ over an
algebraically closed field of arbitrary characteristic, let $L$ be very ample,
and let $m\geq1$.  For a nonempty finite reduced set
$Z\subset X$, suppose that in the complete $L^m$-embedding the vector span of
the points of $Z$ contains the affine embedded tangent space at every point of
$Z$.  Equivalently,
\[
 H^0(X,\cI_Z\otimes L^m)=H^0(X,\cI_{2Z}\otimes L^m).
\]
We prove the exact codimension-free lower bound
\[
 \rk\bigl(H^0(X,L^m)\to H^0(Z,L^m|_Z)\bigr)
 \geq\binom{d+m}{d},\qquad |Z|\geq\binom{d+m}{d}.
\]
The proof projects the original $L$-embedding finitely onto $\PP^d$, injective
on $Z$ and with isomorphic tangent maps there.  Tangent absorption says that
degree-$m$ equations of the projected points survive first fattening.  Over a
perfect field, the inequality
$\alpha(I_Y^{(2)})\geq\alpha(I_Y)+1$ follows by minimal degree: a nonzero
partial derivative lowers the degree, while vanishing of all partials in
characteristic $p$ makes the form a $p$th power, and radicality lowers the
degree again.  The binomial floor is sharp for every pair $(d,m)$ over
$\CC$, including with a proper point span: simplex-lattice interpolation
points are placed in a hyperplane tangent to a smooth hypersurface at every
support.  We also prove
downward propagation to every $L^k$, $k\leq m$, and combine the exact floor
with independent first-jet blocks and a Gauss-fibre alternative.  The
ingredients are classical.  The contribution is their combination into an
exact absorption theorem in arbitrary characteristic; no novelty is claimed
for the derivative/$p$th-root fattening argument itself.
\end{abstract}

\section{Introduction}

Let $\kk$ be an algebraically closed field, let $X^d$ be a smooth projective
integral $\kk$-variety, and let $L$ be a very ample line bundle.  The complete
linear system embeds
\[
 \phi_L:X\hookrightarrow\PP\bigl(H^0(X,L)^*\bigr).
\]
For a finite reduced set $Z\subset X$, write
\[
 S_Z(m):=\Span_{\kk}\{\operatorname{ev}_p:p\in Z\}
          \subset H^0(X,L^m)^*,
 \qquad h_Z(m):=\dim_{\kk}S_Z(m).
\]
In the complete $L^m$-embedding, the affine tangent space at $p$ is the dual
of the first-jet quotient $H^0(2p,L^m|_{2p})$.  We study the rigid incidence
\begin{equation}\label{eq:absorption-intro}
 \widehat T_{\phi_{L^m}(p)}\phi_{L^m}(X)\subseteq S_Z(m)
 \qquad(p\in Z),
\end{equation}
called \emph{point-span tangent absorption}.  It is stronger than an ordinary
Terracini defect: the span of the reduced evaluations already contains all
first-order information at the same supports.

\begin{theorem}[Exact universal floor]\label{thm:exact-floor}
Let $\kk$ be algebraically closed of arbitrary characteristic.  Let $X$ be a
smooth projective integral $\kk$-variety of dimension $d\geq1$, let $L$ be
very ample, let $m\geq1$, and let $Z\subset X$ be a nonempty finite reduced
set.  If $Z$ is point-span tangent-absorbing for the complete system $L^m$,
then
\begin{equation}\label{eq:exact-floor}
 h_Z(m)\geq \BZ:=\binom{d+m}{d},
 \qquad |Z|\geq \BZ.
\end{equation}
\end{theorem}

There is no dependence on the degree, codimension, equations, projective
normality, or Gauss-map degeneracy of the original embedding.  The statement
also includes $m=1$ and $m=2$, where a previous squared-hyperplane argument
did not apply.

The proof has two structural steps.  First choose a finite projection
$\pi:X\to\PP^d$ which is injective on $Z$ and whose tangent map is an
isomorphism at every point of $Z$.
If $Y=\pi(Z)$ and a degree-$m$ form $F$ vanishes on $Y$, its pullback vanishes
on $Z$.  Absorption makes the pullback vanish to second order.  The cotangent
isomorphisms transfer this back to $F$, giving
$(I_Y)_m=(I_Y^{(2)})_m$.  Second, a nonzero form of least degree in the radical
ideal $I_Y$ cannot vanish doubly along $Y$.  In characteristic zero a nonzero
partial derivative lowers the degree.  In characteristic $p$, either the same
happens or every partial vanishes, in which case perfection gives $F=G^p$ and
radicality puts the lower-degree form $G$ in $I_Y$.  Consequently
$\alpha(I_Y^{(2)})\geq\alpha(I_Y)+1$ in every characteristic, and the displayed
equality forces $(I_Y)_m=0$.  The $\binom{d+m}{d}$ degree-$m$ forms on $\PP^d$
then evaluate independently on $Y$.

The floor is not merely an ambient saturation bound.

\begin{theorem}[Proper-span equality in every dimension and degree]
\label{thm:sharpness}
For every $d,m\geq1$ there exist over $\CC$ a smooth integral hypersurface
$X^d\subset\PP^{d+1}$, a nonempty reduced set $Z\subset X$, and
$L=\cO_X(1)$ such that
\[
 |Z|=h_Z(m)=\binom{d+m}{d},\qquad
 h^0(X,L^m)=\binom{d+m+1}{d+1}>h_Z(m),
\]
and $Z$ is point-span tangent-absorbing for $L^m$.  Moreover, every
projective tangent space $T_pX$, $p\in Z$, may be the same hyperplane.
\end{theorem}

Thus neither floor can be raised, even after excluding configurations whose
point span is the entire $L^m$-ambient space.  The equality construction uses
the $\binom{d+m}{d}$ lattice points
\[
 [1:\alpha_1:\cdots:\alpha_d:0],\qquad
 \alpha\in\NN^d,\quad |\alpha|\leq m,
\]
which are unisolvent for polynomials of total degree at most $m$.  A general
high-degree hypersurface with prescribed common tangent hyperplane at these
points is smooth.  On it, the normal coordinate belongs to the square of
every corresponding maximal ideal, converting each degree-$m$ equation of
$Z$ into a double equation.

For $m\geq3$ the exact theorem combines with a rank-sensitive first-jet
argument.  Put
\[
 r_1(Z):=\rk\bigl(H^0(X,L)\to H^0(Z,L|_Z)\bigr)
\]
and let $g_L(Z)$ be the largest number of points of $Z$ in one fibre of the
Gauss map of the $L$-embedding.  We obtain
\begin{equation}\label{eq:combined-intro}
 h_Z(m),|Z|\geq
 \max\left\{\binom{d+m}{d},(d+1)r_1(Z)\right\}.
\end{equation}
Absorption at degree $m$ descends to degree one.  Hence $r_1(Z)\geq d+1$;
if equality holds, every original tangent space along $Z$ is one fixed
$d$-plane.  If $g_L(Z)<\binom{d+m}{d}$, then
\begin{equation}\label{eq:gauss-intro}
 h_Z(m),|Z|\geq
 \max\left\{\binom{d+m}{d},(d+1)(d+2)\right\}.
\end{equation}

\paragraph{Relation to prior work.}
The inequality $\alpha(I^{(2)})\geq\alpha(I)+1$ for finite reduced points is
standard in the study of fattening and symbolic powers.  Bocci--Chiantini
\cite[Remark~2.2]{BocciChiantini2011} give explicitly, for planar point sets
over an algebraically closed field of arbitrary characteristic, the same
derivative/$p$th-root dichotomy used below; that argument is independent of
the projective dimension.  See also Bauer--Szemberg
\cite[Lemma~2.4]{BauerSzemberg2014} for the
arbitrary-projective-space differential inequality in their setting.  The
Zariski--Nagata description of symbolic powers by differential operators over
perfect fields is surveyed by Dao--De Stefani--Grifo--Huneke--N\'u\~nez-
Betancourt \cite[Proposition~2.14 and Exercise~2.15]{DaoEtAl2018}; its
mixed-characteristic extension is developed by De Stefani--Grifo--Jeffries
\cite{DeStefaniGrifoJeffries2020}.  Our degree-two proof is included to make
the positive-characteristic input explicit and self-contained.  Finite general
projections and Bertini theorems with prescribed base schemes are classical;
compare Altman--Kleiman \cite{AltmanKleiman1979} and the analogous adapted-
projection construction over finite fields in Wang \cite[Lemma~A.7]{Wang2022}.
Terracini loci and strong
tangential base loci concern adjacent, but weaker, incidences
\cite{BallicoChiantini2021,BallicoBrambillaSantarsiero2026}.  Jet-ampleness
provides a broader language for infinitesimal interpolation
\cite{BeltramettiDiRoccoSommese1999}.

The novelty claim is deliberately narrow: an adapted finite projection and
the standard fattening inequality are combined to obtain the exact
codimension-free rank floor in arbitrary characteristic, and smooth
proper-span extremizers are given over $\CC$ for every $(d,m)$.  No novelty is
claimed for those ingredients separately or for the positive-characteristic
fattening argument.  The bibliographic search is not a priority certification,
and no peer review is claimed.

\section{Ranks, double points, and absorption}\label{sec:setup}

At a smooth point $p\in X$, write $2p$ for the first infinitesimal
neighbourhood defined by $\mathfrak m_p^2$; it has length $d+1$.  For a
finite reduced set $Z$, put $2Z=\coprod_{p\in Z}2p$.  Since $L^m$ is very
ample, restriction to $2p$ is surjective, and its dual is the affine tangent
space to the cone over the $L^m$-embedding.

\begin{definition}
A nonempty finite reduced set $Z\subset X$ is \emph{point-span
tangent-absorbing for $L^m$} if \eqref{eq:absorption-intro} holds for every
$p\in Z$.
\end{definition}

\begin{lemma}[Kernel criterion]\label{lem:kernel}
The following are equivalent:
\begin{enumerate}[label=(\roman*)]
\item $Z$ is point-span tangent-absorbing for $L^m$;
\item $H^0(X,\cI_Z\otimes L^m)=H^0(X,\cI_{2Z}\otimes L^m)$;
\item restriction to $Z$ and to $2Z$ has the same rank.
\end{enumerate}
\end{lemma}

\begin{proof}
The annihilator of $S_Z(m)$ is the kernel of restriction to $Z$.  The
annihilator of the span of all affine tangent spaces at the supports is the
kernel of restriction to $2Z$.  The point span is contained in the tangent
span, so equality of spans, annihilators, and ranks are equivalent.
\end{proof}

\begin{lemma}[Downward propagation]\label{lem:downward}
If $Z$ is tangent-absorbing for $L^m$, then it is tangent-absorbing for $L^k$
for every $1\leq k\leq m$.
\end{lemma}

\begin{proof}
Fix $s\in H^0(X,\cI_Z\otimes L^k)$ and $p\in Z$.  Since $L^{m-k}$ is
globally generated, choose $t_p\in H^0(X,L^{m-k})$ with $t_p(p)\neq0$; for
$k=m$, take $t_p=1$.  Then $st_p$ vanishes on $Z$, so absorption at degree
$m$ makes it vanish on $2Z$.  In the local ring at $p$, $t_p$ is a unit;
hence $s\in\mathfrak m_p^2L^k_p$.  Repeating for every $p$ gives equality of
the kernels at degree $k$, and Lemma~\ref{lem:kernel} applies.
\end{proof}

\section{The fattening gap over perfect fields}\label{sec:fattening}

For a nonempty finite reduced $Y\subset\PP^d$, let $I_Y$ be its homogeneous
radical ideal and $I_Y^{(2)}=\bigcap_{q\in Y}\mathfrak m_q^2$.  For a
nonzero homogeneous ideal $J$, let $\alpha(J)$ be its least nonzero degree.

\begin{lemma}[Fattening gap]\label{lem:fattening-gap}
Let $\kk$ be a perfect field and let $Y\subset\PP^d_{\kk}$ be a nonempty
finite reduced set of $\kk$-rational points.  Then
\begin{equation}\label{eq:alpha-gap}
 \alpha(I_Y^{(2)})\geq\alpha(I_Y)+1.
\end{equation}
Consequently, if $(I_Y)_m=(I_Y^{(2)})_m$ for some $m\geq1$, then
$(I_Y)_m=0$.
\end{lemma}

\begin{proof}
Put $a=\alpha(I_Y)\geq1$.  If $0\neq F\in(I_Y^{(2)})_a$, every first partial
derivative of $F$ vanishes on $Y$ and hence belongs to the radical ideal
$I_Y$.  If some partial derivative is nonzero, it lies in $(I_Y)_{a-1}$,
contradicting the definition of $a$.

It remains to consider characteristic $p>0$ when all first partials vanish.
Every exponent occurring in $F$ is then divisible by $p$.  Since $\kk$ is
perfect, the coefficients have $p$th roots, so $F=G^p$ for a nonzero
homogeneous form $G$ of degree $a/p<a$.  Radicality of $I_Y$ and
$F\in I_Y$ give $G\in I_Y$, again contradicting minimality.  This proves
\eqref{eq:alpha-gap} in every characteristic.

For the consequence, suppose $(I_Y)_m\neq0$ and choose
$0\neq F\in(I_Y)_a$ with $a=\alpha(I_Y)\leq m$.  For each $q\in Y$, choose
a degree-$(m-a)$ form $H_q$ with $H_q(q)\neq0$.  Then
$H_qF\in(I_Y)_m=(I_Y^{(2)})_m\subset\mathfrak m_q^2$.  Since $H_q$ is a
local unit at $q$, $F\in\mathfrak m_q^2$.  Thus $F\in I_Y^{(2)}$, contrary
to \eqref{eq:alpha-gap} in degree $a$.
\end{proof}

\section{Finite projections adapted to a finite set}\label{sec:projection}

Replacing the ambient space of the $L$-embedding by the linear span of $X$
does not change $L$ or any evaluation rank.

\begin{lemma}[Adapted finite projection]\label{lem:projection}
Let $X^d\subset\PP^M$ be smooth, nondegenerate, projective, and integral,
$d\geq1$, and let $Z\subset X$ be finite and reduced.  There is a linear
projection $\pi_\Lambda:X\to\PP^d$ which is finite, injective on $Z$, has
$d\pi_p:T_pX\xrightarrow{\sim}T_{\pi(p)}\PP^d$ for every $p\in Z$, and satisfies
$\pi_\Lambda^*\cO_{\PP^d}(1)=\cO_X(1)$.  If $M=d$, take the identity after
identifying $X=\PP^d$.
\end{lemma}

\begin{proof}
If $M=d$, a closed integral $d$-fold in $\PP^d$ is all of $\PP^d$.
Assume $M>d$.  Choose a general
$\Lambda\cong\PP^{M-d-1}$ disjoint from $X$, from the finitely many
projective tangent spaces $T_pX$ for $p\in Z$, and from the finitely many
secant lines $\overline{pq}$ for distinct $p,q\in Z$.  Each avoidance
condition is a nonempty open condition on the Grassmannian.  Their finite
intersection is nonempty because the Grassmannian is irreducible.

Projection from $\Lambda$ is a morphism on $X$.  A fibre is the intersection
of $X$ with a linear $\PP^{M-d}$ containing $\Lambda$ as a hyperplane.  A
positive-dimensional projective closed subset of that $\PP^{M-d}$ meets
every hyperplane, so such a fibre would meet $\Lambda$, contradicting
$X\cap\Lambda=\varnothing$.  The morphism is therefore quasi-finite and
projective, hence finite.  Avoiding $T_pX$ makes $d\pi_p$ injective, thus an
isomorphism.  Equivalently, the induced cotangent map is an isomorphism; this
is the only local property used below.  Avoiding secants makes $\pi|_Z$
injective.  The pullback identity follows
from the linear system defining the projection.
\end{proof}

\section{Proof of the exact floor}\label{sec:exact-proof}

\begin{proof}[Proof of Theorem~\ref{thm:exact-floor}]
Embed $X$ by $L$, apply Lemma~\ref{lem:projection}, and put $Y=\pi(Z)$.
The map $Z\to Y$ is a bijection.  If $F\in(I_Y)_m$, then
$s=\pi^*F\in H^0(X,L^m)$ vanishes on $Z$.  Absorption makes $s$ vanish on
$2Z$.  At $p\in Z$, with $q=\pi(p)$,
\[
 d(s)_p=d(F)_q\circ d\pi_p.
\]
The left side is zero and $d\pi_p$ is an isomorphism, so $d(F)_q=0$.
Therefore $F\in\mathfrak m_q^2$ for all $q\in Y$ and
$(I_Y)_m=(I_Y^{(2)})_m$.  Lemma~\ref{lem:fattening-gap} gives
$(I_Y)_m=0$.

Evaluation on $Y$ is consequently injective on
$H^0(\PP^d,\cO(m))$, of dimension $\binom{d+m}{d}$.  Pulling this subspace
back to $H^0(X,L^m)$ and using $Z\simeq Y$ proves
$h_Z(m)\geq\binom{d+m}{d}$.  Since $S_Z(m)$ is spanned by $|Z|$ evaluation
vectors, $|Z|\geq h_Z(m)$.
\end{proof}

\begin{corollary}[Projective form]
Under the hypotheses of Theorem~\ref{thm:exact-floor},
$\dim\Span_{\PP}\phi_{L^m}(Z)\geq\binom{d+m}{d}-1$.
\end{corollary}

\begin{corollary}[No low-cardinality extreme defect]
If $|Z|<\binom{d+m}{d}$, thickening $Z$ to $2Z$ strictly increases the rank
of restriction in $L^m$.
\end{corollary}

\begin{remark}
Perfection and radicality are the inputs used by this positive-characteristic
proof.  The theorem is stated over an algebraically closed field, hence over a
perfect field, and the finite support is reduced.  No analogous conclusion is
asserted here for imperfect ground fields or nonreduced supports.
\end{remark}

\section{Simplex-lattice unisolvence}\label{sec:unisolvence}

Here $\NN=\{0,1,2,\ldots\}$.

Fix $d,m\geq1$ and set
\[
 A_{d,m}:=\{\alpha=(\alpha_1,\ldots,\alpha_d)\in\NN^d:|\alpha|\leq m\},
 \qquad |A_{d,m}|=\binom{d+m}{d}=\BZ.
\]
For $\alpha\in A_{d,m}$, let
$p_\alpha=[1:\alpha_1:\cdots:\alpha_d]\in\PP^d$.

\begin{lemma}[Simplex-lattice interpolation]\label{lem:unisolvent}
Evaluation at the points $p_\alpha$, $\alpha\in A_{d,m}$, is an isomorphism
\[
 H^0(\PP^d,\cO_{\PP^d}(m))\xrightarrow{\ \sim\ }\CC^{A_{d,m}}.
\]
\end{lemma}

\begin{proof}
On the chart $x_0=1$, consider the falling-factorial polynomials
\[
 N_\beta(t_1,\ldots,t_d):=\prod_{i=1}^d\binom{t_i}{\beta_i},
 \qquad \beta\in A_{d,m}.
\]
They form a basis of the polynomials of total degree at most $m$, since each
has leading monomial $\prod_i t_i^{\beta_i}/\beta_i!$.  Order $A_{d,m}$ by a
linear extension of the coordinatewise partial order.  Then
$N_\beta(\alpha)=0$ unless $\beta_i\leq\alpha_i$ for every $i$, while
$N_\alpha(\alpha)=1$.  The evaluation matrix is triangular with diagonal
one.  Homogenization proves the assertion.
\end{proof}

This is a concrete multivariate Newton interpolation set.  The proof is
included so the equality construction does not depend on numerical
conditioning or genericity; related lattice constructions appear in
\cite{ChungYao1977,LeePhillips1991}.

\section{Smooth common-tangent extremizers}\label{sec:extremizers}

Throughout this section the ground field is $\CC$.

Let $W=V(y)\cong\PP^d$ be a hyperplane in $\PP^{d+1}$ with coordinates
$[x_0:\cdots:x_d:y]$.  Regard the simplex-lattice set
\[
 Z=\{[1:\alpha_1:\cdots:\alpha_d:0]:\alpha\in A_{d,m}\}\subset W.
\]
Write $2Z_W$ for the union of first neighbourhoods inside $W$.

\begin{proposition}[Smooth prescribed-tangent hypersurface]
\label{prop:bertini}
Put $r=|Z|=\BZ$.  For every $n\geq2r+1$, there is a smooth integral
degree-$n$ hypersurface $X\subset\PP^{d+1}$ containing $Z$ such that
$T_pX=W$ for every $p\in Z$.
\end{proposition}

\begin{proof}
Consider forms
\begin{equation}\label{eq:F-family}
 F(x,y)=f(x_0,\ldots,x_d)+yG(x_0,\ldots,x_d,y),
\end{equation}
where $f\in H^0(W,\cI_{2Z_W}(n))$ and
$G\in H^0(\PP^{d+1},\cO(n-1))$.

The set-theoretic support of this linear system's base scheme is exactly
$Z$; no assertion about its scheme structure is needed.  Outside $W$, the
term $yG$ is base-point-free.  If $q\in W\setminus Z$, for each $p\in Z$ choose a
hyperplane $\ell_p\subset W$ through $p$ but not through $q$.  Then
$\prod_{p\in Z}\ell_p^2$ has degree $2r$, vanishes doubly at $Z$, and is
nonzero at $q$; multiply it by a degree-$(n-2r)$ form nonzero at $q$.

Thus the restricted linear system on the smooth open set
$\PP^{d+1}\setminus Z$ is base-point-free.  The characteristic-zero smooth
Bertini theorem gives smoothness for a general member on that open set.  At
$p\in Z$, all tangential derivatives of $f$ vanish, whereas the normal
derivative of $F$ is $G(p)$.  The finitely many conditions $G(p)\neq0$ define
a nonempty open set.  Hence a general member is smooth at $Z$ and has tangent
hyperplane $W$ there.  It is integral: if a projective hypersurface of
positive dimension were reducible, two positive-degree components would
meet, and their union would be singular along the intersection.  Bertini
with a prescribed base scheme is
standard; see Altman--Kleiman \cite{AltmanKleiman1979}.
\end{proof}

The numerical threshold is a transparent sufficient bound; no minimality in
$n$ is claimed.

\begin{proof}[Proof of Theorem~\ref{thm:sharpness}]
Take $Z\subset W$ as above and choose $X=V(F)$ from
Proposition~\ref{prop:bertini}, with $n\geq2\BZ+1$.  Since $n>m$ and
$d+1\geq2$, the hypersurface sequence
\[
0\longrightarrow\cO_{\PP^{d+1}}(m-n)
\longrightarrow\cO_{\PP^{d+1}}(m)
\longrightarrow\cO_X(m)\longrightarrow0
\]
together with $H^0(\cO(m-n))=0$ and
$H^1(\PP^{d+1},\cO(m-n))=0$ gives
\begin{equation}\label{eq:restriction-iso}
 H^0(\PP^{d+1},\cO(m))\xrightarrow{\ \sim\ }H^0(X,\cO_X(m)).
\end{equation}

Let a degree-$m$ form $Q(x,y)$ vanish on $Z$.  Its restriction $Q(x,0)$ to
$W$ vanishes at the unisolvent set of Lemma~\ref{lem:unisolvent}, hence is
zero.  Thus $Q=yR$ for a degree-$(m-1)$ form $R$.

Fix $p\in Z$.  Locally, \eqref{eq:F-family} reads $f+yG=0$, where
$f\in\mathfrak m_{W,p}^2$ and $G(p)\neq0$.  Thus $G$ is a unit and, in
$\cO_{X,p}$,
\[
 y=-f/G\in\mathfrak m_{X,p}^2.
\]
Therefore $Q=yR$ vanishes on $2p$.  This holds for every $p$, proving
\[
 H^0(X,\cI_Z\otimes\cO_X(m))
 =H^0(X,\cI_{2Z}\otimes\cO_X(m)).
\]
The set is absorbing by Lemma~\ref{lem:kernel}.

Lemma~\ref{lem:unisolvent} shows that the evaluations have rank
$|Z|=\BZ$.  Equation \eqref{eq:restriction-iso} gives
$h^0(X,\cO_X(m))=\binom{d+m+1}{d+1}>\BZ$, so the point span is proper.  The
common-tangent assertion is part of Proposition~\ref{prop:bertini}.
\end{proof}

\begin{remark}
The equality family explains why no improvement follows merely from
excluding ambient point spans.  Any stronger theorem needs geometric input
that rules out, controls, or classifies high-contact common-tangent sets.
\end{remark}

\section{Rank-sensitive jet blocks and Gauss fibres}\label{sec:rank-sensitive}

Return to an arbitrary smooth $X$ and very ample $L$, and put
$r_1(Z)=h_Z(1)$.  The following extension is useful for $m\geq3$.

\begin{lemma}[Squared-hyperplane extension]\label{lem:block-extension}
Let $m\geq3$.  Suppose $p_1,\ldots,p_s\in X$ have independent first
neighbourhoods for $L^m$.  If the $L$-evaluation at $x\in X$ is not in the
span of the evaluations at the $p_i$, then
$2p_1\cup\cdots\cup2p_s\cup2x$ also imposes independent conditions on
$H^0(X,L^m)$.
\end{lemma}

\begin{proof}
In the $L$-embedding choose $e\in H^0(X,L)$ vanishing at every $p_i$ but not
at $x$.  The section $e^2$ vanishes on the old double points and is a unit on
$2x$.  Since $L^{m-2}$ is very ample, restriction to $2x$ is surjective.
Thus $e^2H^0(X,L^{m-2})$ realizes arbitrary first-jet data at $x$ while
vanishing on the old union.  Correcting residual data at $x$ proves the claim.
\end{proof}

\begin{proposition}[Rank-sensitive floor]\label{prop:rank-sensitive}
If $m\geq3$ and $Z$ is tangent-absorbing for $L^m$, then
\[
 h_Z(m),|Z|\geq
 \max\left\{\binom{d+m}{d},(d+1)r_1(Z)\right\}.
\]
\end{proposition}

\begin{proof}
Choose $r_1(Z)$ points whose degree-one evaluations form a basis, ordered so
that each new evaluation escapes the preceding span.  A single double point
imposes $d+1$ independent conditions.  Repeated application of
Lemma~\ref{lem:block-extension} makes the $r_1(Z)$ double points independent
in $L^m$.  By absorption, their $(d+1)r_1(Z)$-dimensional dual first-jet span
lies in $S_Z(m)$.  Combine this with Theorem~\ref{thm:exact-floor} and
$|Z|\geq h_Z(m)$.
\end{proof}

Let
\[
 \gamma_L:X\longrightarrow\Gr(d+1,H^0(X,L)^*)
\]
be the Gauss map of the $L$-embedding and define
\[
 g_L(Z):=\max_{W\in\Gr(d+1,H^0(X,L)^*)}
 |Z\cap\gamma_L^{-1}(W)|.
\]

\begin{proposition}[Degree-one alternative]\label{prop:degree-one}
Under Theorem~\ref{thm:exact-floor}, $r_1(Z)\geq d+1$.  If
$r_1(Z)=d+1$, the degree-one evaluation span is the affine tangent space at
every point of $Z$; all of $Z$ lies in one fibre of $\gamma_L$ and in its
projective tangent $d$-plane.
\end{proposition}

\begin{proof}
By Lemma~\ref{lem:downward}, $Z$ is absorbing for $L$.  Its degree-one
evaluation span contains a $(d+1)$-dimensional affine tangent space.  If its
dimension is exactly $d+1$, every such tangent space along $Z$ is a subspace
of the same dimension inside it, hence equals it.  Evaluation representatives
also lie in their own affine tangent spaces.
\end{proof}

\begin{corollary}[Gauss-fibre refinement]\label{cor:gauss}
Let $m\geq3$.  If $Z$ is absorbing for $L^m$ and
$g_L(Z)<\binom{d+m}{d}$, then
\[
 h_Z(m),|Z|\geq
 \max\left\{\binom{d+m}{d},(d+1)(d+2)\right\}.
\]
\end{corollary}

\begin{proof}
If $r_1(Z)=d+1$, Proposition~\ref{prop:degree-one} puts all $Z$ in one Gauss
fibre, while Theorem~\ref{thm:exact-floor} gives
$|Z|\geq\binom{d+m}{d}$, a contradiction.  Hence $r_1(Z)\geq d+2$, and
Proposition~\ref{prop:rank-sensitive} applies.
\end{proof}

For a smooth quadric, the Gauss map is injective on geometric points, so the
corollary applies automatically.  Indeed, if two points have the same polar
hyperplane, after rescaling their difference lies in the radical of the polar
bilinear form.  If that difference were nonzero, the quadratic identity and
the fact that both points lie on the quadric would make it a singular point of
the quadric.  This also covers characteristic two, where the Gauss map may be
inseparable and ``polarity'' alone would be imprecise.  When the binomial term
already dominates $(d+1)(d+2)$, the exact theorem is stronger.

\section{Comparison with adjacent interpolation loci}\label{sec:comparison}

\begin{center}
\begin{tabular}{p{0.25\textwidth}p{0.64\textwidth}}
\toprule
Object & Incidence condition \\
\midrule
Terracini locus & The union $2Z$ fails to impose the expected independent
conditions; the tangent span is defective. \\
Strong tangential base locus & A span generated by selected tangent spaces
contains another point or tangent space. \\
Point-span tangent absorption & The span of reduced evaluations already
contains the tangent space at every support; point rank equals double-point
rank. \\
Equality family here & One hyperplane is tangent at an unisolvent set and its
normal coordinate is second order in each local ring, giving exact binomial
rank with proper point span. \\
\bottomrule
\end{tabular}
\end{center}

Absorption implies an extreme Terracini defect unless the point span is
ambient, but ordinary Terracini defect does not imply absorption.  Nor can
one generally delete a support and reinterpret absorption as a strong-base-
locus condition: its evaluation direction may be needed to express its own
tangent space.  Theorem~\ref{thm:exact-floor} bypasses those classification
problems by transporting the incidence to points in $\PP^d$ and using initial
degree.

\section{Exact computational witnesses}\label{sec:replay}

The accompanying replay
\texttt{repro/verify\_exact\_projection\_floor.py} uses exact rational
arithmetic to construct simplex-lattice evaluation matrices for a grid of
$(d,m)$.  It verifies rank $\binom{d+m}{d}$ and checks a falsification
boundary: deleting one node leaves a nonzero degree-$m$ equation, while
adjoining first-jet rows increases the value rank.

The replay
\texttt{repro/verify\_common\_tangent\_extremizer.py} checks explicit local
fixtures of the family $F=f+yG$.  It verifies unisolvence, prescribed normal
gradients at the supports, the local second-order identity for $y$, and
equality of value and double-point ranks.  These are exact finite fixtures.
They do not replace Bertini, certify every member, prove the universal
theorem, establish priority, or constitute formal verification.

The positive-characteristic replay
\texttt{repro/verify\_perfect\_field\_fattening.py} works by exact Gaussian
elimination modulo $p$.  It exhausts every nonempty reduced support in
$\PP^1(\mathbf F_2)$, $\PP^1(\mathbf F_3)$, and
$\PP^2(\mathbf F_2)$ through explicit stated degree cutoffs.  For each
support it checks the initial-degree gap and, degree by degree, that equality
of the value and first-neighbourhood kernels occurs only when the value kernel
is zero.  Separate fixtures exhibit the all-partials-zero case $F=G^p$ and
the lower-degree radical root $G$.  Its exhaustiveness is confined to these
finite spaces and degree ranges.

A counterexample to the main theorem would be a smooth $d$-fold over an
algebraically closed field, a very ample $L$, and a nonempty reduced absorbing $Z$ with
$h_Z(m)<\binom{d+m}{d}$.  After an adapted projection it would either violate
the transferred equality $(I_Y)_m=(I_Y^{(2)})_m$ or the perfect-field
fattening gap.

\section{Scope, limitations, and provenance}\label{sec:scope}

\begin{itemize}[leftmargin=*]
\item The exact floor is proved over an algebraically closed field of arbitrary
characteristic.  The proof of the fattening gap uses perfection; no descent to
imperfect fields is claimed.
\item $X$ is smooth, projective, integral, and pure-dimensional, and $Z$ is
nonempty, finite, and reduced.  Nonreduced supports, singular ambient schemes,
and mixed dimensions need different local statements.
\item The proper-span sharpness construction is retained over $\CC$ and proves
existence for a sufficient hypersurface degree.  It neither supplies a
positive-characteristic Bertini realization, classifies equality, nor
minimizes that degree.
\item The Gauss-fibre result is a secondary refinement for $m\geq3$.  It does
not classify Gauss fibres, contact loci, Terracini loci, or strong base loci.
\item The work proves no statement about a Millennium Prize Problem.  The
0--10 research-impact scores used during drafting are editorial heuristics,
not mathematical claims.
\item Literature searches covered the cited fat-point, projection, Bertini,
interpolation, Terracini, jet-ampleness, and Gauss-map sources.  They are not
exhaustive MathSciNet/zbMATH review or specialist peer review.
\item The replay scripts are finite diagnostic witnesses.  No Lean
certification, independent reproduction, or external validation is claimed.
\end{itemize}

The extension of the exact absorption theorem from characteristic zero to
arbitrary characteristic applies the standard Frobenius--radical alternative
in the minimal-degree fattening lemma.  The proper-span complex extremizers
and secondary block refinements are retained.
OpenAI Codex assisted with proof discovery and auditing, literature triage,
replay code, and drafting.  The named author is responsible for the claims
and any public submission.  No external human peer review or priority
certification is claimed.

\small
\begin{thebibliography}{99}

\bibitem{AltmanKleiman1979}
A.~B.~Altman and S.~L.~Kleiman,
\emph{Bertini theorems for hypersurface sections containing a subscheme},
Comm. Algebra 7 (1979), no.~8, 775--790.
\url{https://doi.org/10.1080/00927877908822375}

\bibitem{BallicoBrambillaSantarsiero2026}
E.~Ballico, M.~C.~Brambilla, and P.~Santarsiero,
\emph{On the strong base locus of a projective variety},
arXiv:2603.15103 (2026).
\url{https://arxiv.org/abs/2603.15103}

\bibitem{BallicoChiantini2021}
E.~Ballico and L.~Chiantini,
\emph{On the Terracini locus of projective varieties},
Milan J. Math. 89 (2021), 1--17.
\url{https://doi.org/10.1007/s00032-020-00324-5}

\bibitem{BauerSzemberg2014}
T.~Bauer and T.~Szemberg,
\emph{The effect of points fattening in dimension three},
arXiv:1308.4983 (2014).
\url{https://arxiv.org/abs/1308.4983}

\bibitem{BeltramettiDiRoccoSommese1999}
M.~C.~Beltrametti, S.~Di Rocco, and A.~J.~Sommese,
\emph{On generation of jets for vector bundles},
Rev. Mat. Complut. 12 (1999), no.~1, 27--45.
\url{https://doi.org/10.5209/rev_REMA.1999.v12.n1.17182}

\bibitem{BocciChiantini2011}
C.~Bocci and L.~Chiantini,
\emph{The effect of points fattening on postulation},
J. Pure Appl. Algebra 215 (2011), no.~1, 89--98.
\url{https://doi.org/10.1016/j.jpaa.2010.04.010}

\bibitem{ChungYao1977}
K.~C.~Chung and T.~H.~Yao,
\emph{On lattices admitting unique Lagrange interpolations},
SIAM J. Numer. Anal. 14 (1977), no.~4, 735--743.
\url{https://doi.org/10.1137/0714050}

\bibitem{DeStefaniGrifoJeffries2020}
A.~De Stefani, E.~Grifo, and J.~Jeffries,
\emph{A Zariski--Nagata theorem for smooth $\mathbf Z$-algebras},
J. Reine Angew. Math. 761 (2020), 123--140.
\url{https://doi.org/10.1515/crelle-2018-0012}

\bibitem{DaoEtAl2018}
H.~Dao, A.~De Stefani, E.~Grifo, C.~Huneke, and L.~N\'u\~nez-Betancourt,
\emph{Symbolic powers of ideals}, in \emph{Singularities and Foliations:
Geometry, Topology and Applications}, Springer Proc. Math. Stat. 222 (2018),
387--432.
\url{https://doi.org/10.1007/978-3-319-73639-6_13}

\bibitem{LeePhillips1991}
S.~L.~Lee and G.~M.~Phillips,
\emph{Construction of lattices for Lagrange interpolation in projective space},
Constr. Approx. 7 (1991), no.~3, 283--297.
\url{https://doi.org/10.1007/BF01888158}

\bibitem{Wang2022}
X.~Wang,
\emph{On the Bertini regularity theorem for arithmetic varieties},
J. Ec. polytech. Math. 9 (2022), 601--670; see Lemma~A.7.
\url{https://doi.org/10.5802/jep.191}

\end{thebibliography}

\end{document}
